
\begin{frame}{Pattern 1: Annual Cycle Limits Are Ubiquitous}
\textbf{Every battery in the corpus has an annual equivalent full cycle cap.}
\vspace{0.3cm}
\begin{itemize}
    \item Range: 250--365 EFCs/year across the 16 projects
    \item 365 EFCs/year is the most common limit---equivalent to one full discharge per day on average
    \item At 4-hour duration, 365 EFCs/year $\approx$ 1,460 MWh of annual throughput per MW of capacity
    \pause
    \item \textbf{Key implication}: cycle budgets create \textit{intertemporal linkages} across dispatch decisions
    \begin{itemize}
        \item An operator facing a binding annual cap treats a cycle today as reducing the option value of cycling in future high-price periods
        \item This can distort dispatch even on days where the daily limit is not binding
    \end{itemize}
    \pause
    \item Some filings also impose \textbf{daily cycle limits} (2 EFCs/day) or refer daily limits to undisclosed manufacturer warranty---creating a gap in the public record
\end{itemize}
\end{frame}

\begin{frame}{Pattern 2: SOC Constraints Vary in Direction and Design}
SOC constraints appear in nearly every ESA, but their design differs substantially:
\vspace{0.2cm}
\begin{itemize}
    \item \textbf{Average SOC ceilings} (most common): battery must not be held too fully charged on average
    \begin{itemize}
        \item Range: 40\%--60\% across the corpus
        \item Penalty design varies: direct financial damages vs.\ retroactive cycle-count reduction
    \end{itemize}
    \pause
    \item \textbf{Average SOC floors} (TX-59034): battery must not be held too discharged on average ($>51\%$)
    \begin{itemize}
        \item Opposite direction---prevents chronic deep-discharge, not chronic overcharge
    \end{itemize}
    \pause
    \item \textbf{Resting SOC limits}: separate ceiling on SOC during prolonged non-operating periods (e.g.\ $<$20\% in TX-57760; $>$30\% in Quail Ranch during outages $>$3 months)
    \pause
    \item \textbf{Minimum hold rule} (NM-25-00048): SOC may not remain below 3\% for more than 1 hour continuously
\end{itemize}
\vspace{0.2cm}
The diversity of SOC constraint design across a small corpus suggests these terms are \textit{manufacturer-specific}, not standardized.
\end{frame}

\begin{frame}{Pattern 3: The Warranty Chain}
Restrictions originate from battery manufacturers but reach grid dispatch through a layered contractual chain:

\vspace{0.3cm}
\begin{center}
\textbf{Manufacturer warranty}\\
$\downarrow$\\
ESA Exhibit L / Exhibit O (``ESS Operating Restrictions'')\\
$\downarrow$\\
Utility dispatch authority (PNM, SPS, EPE)\\
$\downarrow$\\
\textit{Regulatory approval} (NMPRC, PUCT)
\end{center}

\vspace{0.3cm}
\begin{itemize}
    \item The Special Service Contract (SSC) between PNM and Greater Kudu LLC (Meta) states: \textit{``In charging and discharging SSC Storage Resources, PNM shall abide by any warranty restrictions or other terms of the Third Party ESAs''} --- \textcolor{gray}{SSC \S8.4.3, NM-25-00048-UT}
    \item \textbf{The NMPRC approved this arrangement}: the Commission is effectively ratifying unnamed manufacturer warranty terms as conditions on utility dispatch
\end{itemize}
\end{frame}

\begin{frame}{Pattern 4: Enforcement Design Is Evolving}
\textbf{Earlier ESAs (2021--2023)}: financial disincentives applied retroactively.
\begin{itemize}
    \item Exceeding the SOC ceiling reduces cycle allowance for that year (NM-21-00031-UT, NM-23-00353-UT)
    \item Violating RTE guarantee triggers payment adjustment (TX-57760)
    \item Operator is not blocked in real time---can exceed the constraint and face financial consequences later
\end{itemize}

\vspace{0.3cm}
\pause
\textbf{Newer ESAs (2025)}: automated real-time refusal.
\begin{itemize}
    \item NM-25-00048-UT (Four Mile Mesa, Star Light, Windy Lane): \textit{``Seller will have no obligation to dispatch the ESS in response to any dispatch instructions provided by PNM that are not in accordance with the Operating Parameters''}
    \item The BESS control system \textbf{physically refuses} the dispatch command
    \item Non-conforming period counts as ``Seller Excused Hours''---no penalty to Seller
    \item Operator \textit{cannot} exceed the constraint; it is enforced at the hardware/software level
\end{itemize}
\vspace{0.2cm}
\pause
\textbf{Implication}: newer batteries may show \textit{hard kinks} in dispatch behavior at constraint boundaries rather than the smooth forward-looking adjustment predicted for soft penalties.
\end{frame}
